% Work/MAT | 1. rocnik | nova hodina 33: Sdruzene vyrazy.
% Navaznost: 31 castecnne odmocnovani, 32 jednoduche usmernovani.
% Puvodni obsah hodiny 36 se presouva sem; hodina 34 je na slozene vyrazy.
% Ucebnice: Mocniny a Odmocniny.tex, Usmernovani zlomku: sdruzeny vyraz.
% Realisticky.cz 020714 Pocitani s odmocninami III: usmernovani dvojclenu,
% prakticke usporadani postupu. Priklady jsou nove sestaveny.
\documentclass[11pt,a4paper]{article}
% Spolecny styl pro ucitelske pripravy MAT (Work/MAT)
% Verze 1.1, 2026-10-23
\NeedsTeXFormat{LaTeX2e}
\ProvidesPackage{priprava}[2026/10/23 v1.1 Sablona ucitelskych priprav MAT]

% Zakladni typografie a cestina
\RequirePackage[T1]{fontenc}
\RequirePackage[utf8]{inputenc}
\RequirePackage[czech]{babel}
\RequirePackage{lmodern}
\RequirePackage{microtype}
\RequirePackage{amsmath,amssymb}
\RequirePackage{xcolor}
\RequirePackage{enumitem}
\RequirePackage[most]{tcolorbox}
\RequirePackage{tabularx,booktabs,array}
\RequirePackage{fancyhdr}
\RequirePackage{needspace}
\RequirePackage{geometry}
\RequirePackage[hidelinks]{hyperref}

\geometry{a4paper,top=17mm,bottom=18mm,left=23mm,right=23mm}
\definecolor{notegray}{gray}{0.39}
\setlength{\parindent}{0pt}
\setlength{\parskip}{2.5pt plus 1pt minus 0.5pt}
\setlist[itemize]{topsep=2pt,itemsep=1pt,parsep=0pt,leftmargin=1.45em}
\setlist[enumerate]{topsep=2pt,itemsep=2pt,parsep=0pt,leftmargin=1.65em}
\setcounter{secnumdepth}{0}
\renewcommand{\arraystretch}{1.12}

% Zahlavi: vlevo rocnik | globalni tema, vpravo cislo: nazev hodiny
% Zapati zustava zcela prazdne.
\pagestyle{fancy}
\fancyhf{}
\renewcommand{\headrulewidth}{0.35pt}
\renewcommand{\footrulewidth}{0pt}
\setlength{\headheight}{14pt}
\newcommand{\nastavPripravu}[4]{%
  \fancyhead[L]{\small #1\enspace\textbar\enspace #2}%
  \fancyhead[R]{\small #3: #4}%
}

% Minimalisticke cernobile boxy. Kazda definice zvlast.
\tcbset{commonlesson/.style={enhanced,breakable=false,colback=white,colframe=black,
  colbacktitle=white,coltitle=black,boxrule=0.45pt,
  arc=0.8mm,outer arc=0.8mm,left=3mm,right=3mm,top=1.6mm,bottom=1.6mm,
  toptitle=1mm,bottomtitle=0.7mm,fonttitle=\bfseries,
  attach boxed title to top left={xshift=2.5mm,yshift=-1.15mm},
  boxed title style={colback=white,colframe=white,boxrule=0pt,left=1mm,right=1mm}}}
\newtcolorbox{zadaniBox}[1][]{commonlesson,title={Zadání},#1}
\newtcolorbox{prikladBox}[1][]{commonlesson,title={Příklad},#1}
\newtcolorbox{definiceBox}[1][]{commonlesson,title={Definice},#1}

% Spolecny box pro zadani a jeho reseni, bez mezery mezi castmi.
% Pouziti:
%   \begin{prikladReseniBox}
%     Zadani ulohy.
%     \tcblower
%     \begin{align*} ... \end{align*}
%   \end{prikladReseniBox}
% Spodni cast ma jemny sedy podklad a delici linku; lze v ni sazet
% vice radku matematiky. Obe casti maji jen jeden spolecny ramecek.
\tcbset{solutionlesson/.style={
  commonlesson,
  enhanced jigsaw,
  breakable,
  colbacklower=black!3,
  segmentation style={solid,draw=black!25,line width=0.35pt},
  fontlower=\small\color{notegray},
  before lower={\noindent\textbf{Řešení:}\enspace},
  leftlower=3mm,rightlower=3mm
}}
\newtcolorbox{prikladReseniBox}[1][]{solutionlesson,title={Příklad},#1}
\newtcolorbox{zadaniReseniBox}[1][]{solutionlesson,title={Zadání},#1}

% Pomucky pro autora pripravy: sedive poznamky; \reseni pro starsi pripravy
\newcommand{\poznamka}[1]{\par\begingroup\small\color{notegray}\itshape
  \textnormal{\textbf{Poznámka pro učitele.}}\enspace #1\par\endgroup}
\newcommand{\reseni}[1]{\par\begingroup\small\color{notegray}
  \textbf{Očekávané řešení:}\enspace #1\par\endgroup}
\newcommand{\blok}[1]{\needspace{5\baselineskip}\subsection*{#1}}
\newcommand{\vyrok}[1]{\uv{#1}}

\nastavPripravu{1. ročník}{Mocniny a odmocniny}{33}{Sdružené výrazy}
\begin{document}
\poznamka{\textbf{45minutové jádro:} 5 min zopakování usměrňování;
8 min rozdíl čtverců a Příklad 1; 8 min Samostatná práce 1 a)--d);
8 min Příklad 2 a Samostatná práce 2 a)--c);
8 min Příklad 3 a práce se znaménkem; 5 min samostatná kontrola;
3 min shrnutí a domácí práce. Proměnné a dvoustupňové úlohy jsou \textbf{rezerva}.
Dnes jde o jediný nový trik: \emph{dva sdružené dvojčleny vytvoří rozdíl čtverců}.}

\blok{Proč nestačí násobit stejnou odmocninou?}
V hodině 32 jsme usměrnili $1/\sqrt2$ násobením $\sqrt2/\sqrt2$.
Co ale se zlomkem $1/(2+\sqrt3)$?
\[
\frac1{2+\sqrt3}\cdot\frac{\sqrt3}{\sqrt3}
=\frac{\sqrt3}{2\sqrt3+3},
\]
takže odmocnina stále zůstává. Potřebujeme v čitateli a jmenovateli
přidat takový násobek, který \emph{odmocninový člen zruší}.

\begin{definiceBox}[title={Definice: sdružené výrazy}]
\textbf{Sdružené výrazy} jsou dvojčleny $u+v$ a $u-v$,
které se liší pouze znaménkem mezi členy.
Jejich součin je rozdíl čtverců:
\[
\boxed{(u+v)(u-v)=u^2-v^2.}
\]
Toto pravidlo platí pro všechna reálná $u,v$.
\end{definiceBox}
\textbf{Proč?} Po roznásobení $u^2-uv+uv-v^2$ se prostřední členy
odečtou. Je-li ve jmenovateli $u+\sqrt a$, násobíme vhodným
$u-\sqrt a$ (a naopak), pokud je tento násobek nenulový.

\begin{prikladReseniBox}[breakable=false,title={Příklad 1: první dvojčlen ve jmenovateli}]
Usměrni $\dfrac1{2+\sqrt3}$ a $\dfrac3{\sqrt5+2}$.
\tcblower
\[
\frac1{2+\sqrt3}\cdot\frac{2-\sqrt3}{2-\sqrt3}
=\frac{2-\sqrt3}{4-3}=\boxed{2-\sqrt3}.
\]
\[
\frac3{\sqrt5+2}\cdot\frac{\sqrt5-2}{\sqrt5-2}
=\frac{3(\sqrt5-2)}{5-4}=\boxed{3\sqrt5-6}.
\]
Ve druhém příkladu je činitel $\sqrt5-2\ne0$ a původní jmenovatel také není nula.
\end{prikladReseniBox}

\begin{zadaniReseniBox}[breakable=false,title={Samostatná práce 1: jeden sdružený výraz stačí}]
Usměrni:
\begin{enumerate}[label=\alph*)]
\item $\dfrac1{\sqrt3+1}$;
\item $\dfrac2{3-\sqrt7}$;
\item $\dfrac4{\sqrt6-\sqrt2}$;
\item $\dfrac5{\sqrt7+\sqrt2}$;
\item $\dfrac3{2+\sqrt2}$;
\item $\dfrac{\sqrt5}{\sqrt5-1}$.
\end{enumerate}
\tcblower
\begin{enumerate}[label=\alph*)]
\item $\dfrac{\sqrt3-1}{3-1}=\dfrac{\sqrt3-1}2$;
\item $\dfrac{2(3+\sqrt7)}{9-7}=3+\sqrt7$;
\item $\dfrac{4(\sqrt6+\sqrt2)}{6-2}=\sqrt6+\sqrt2$;
\item $\dfrac{5(\sqrt7-\sqrt2)}{7-2}=\sqrt7-\sqrt2$;
\item $\dfrac{3(2-\sqrt2)}{4-2}=\dfrac{3(2-\sqrt2)}2$;
\item $\dfrac{\sqrt5(\sqrt5+1)}{5-1}=\dfrac{5+\sqrt5}{4}$.
\end{enumerate}
\end{zadaniReseniBox}

\blok{V čitateli je také odmocnina: nejprve udrž přehled}
\begin{prikladReseniBox}[breakable=false,title={Příklad 2: usměrňování s dvojčlenem v čitateli}]
Usměrni $\dfrac{\sqrt3+\sqrt2}{\sqrt3-\sqrt2}$.
\tcblower
\[
\frac{\sqrt3+\sqrt2}{\sqrt3-\sqrt2}
\cdot\frac{\sqrt3+\sqrt2}{\sqrt3+\sqrt2}
=\frac{(\sqrt3+\sqrt2)^2}{3-2}
=\boxed{5+2\sqrt6}.
\]
\textbf{Pozor:} $(u+v)^2=u^2+2uv+v^2$, nikoliv pouze $u^2+v^2$.
\end{prikladReseniBox}

\begin{zadaniReseniBox}[breakable=false,title={Samostatná práce 2: promysli nejprve strategii}]
\begin{enumerate}[label=\alph*)]
\item $\dfrac{\sqrt2+1}{\sqrt2-1}$;
\item $\dfrac{\sqrt7-\sqrt3}{\sqrt7+\sqrt3}$;
\item $\dfrac2{\sqrt7-\sqrt3}$;
\item $\dfrac{\sqrt5-1}{\sqrt5+1}$;
\item $\dfrac{\sqrt6+\sqrt2}{\sqrt6-\sqrt2}$.
\end{enumerate}
\tcblower
\begin{enumerate}[label=\alph*)]
\item $(\sqrt2+1)^2/(2-1)=3+2\sqrt2$;
\item $(\sqrt7-\sqrt3)^2/(7-3)=(10-2\sqrt{21})/4=(5-\sqrt{21})/2$;
\item $2(\sqrt7+\sqrt3)/(7-3)=(\sqrt7+\sqrt3)/2$;
\item $(\sqrt5-1)^2/(5-1)=(6-2\sqrt5)/4=(3-\sqrt5)/2$;
\item $(\sqrt6+\sqrt2)^2/(6-2)=(8+4\sqrt3)/4=2+\sqrt3$.
\end{enumerate}
\end{zadaniReseniBox}

\blok{Časté chyby: výběr znaménka, hranice použití}
\begin{prikladReseniBox}[breakable=false,title={Příklad 3: najdi a oprav chybu}]
Žák napíše $\dfrac1{3-\sqrt5}\cdot\dfrac{3-\sqrt5}{3-\sqrt5}
=\dfrac{3-\sqrt5}{9-5}$. Kde je chyba?
\tcblower
Použil \emph{stejný}, nikoliv \emph{sdružený} výraz;
$(3-\sqrt5)^2\ne9-5$.
Správně:
\[
\frac1{3-\sqrt5}\cdot\frac{3+\sqrt5}{3+\sqrt5}
=\boxed{\frac{3+\sqrt5}{4}}.
\]
\end{prikladReseniBox}

\begin{zadaniReseniBox}[breakable=false,title={Samostatná práce 3: tři otázky ke kontrole}]
\begin{enumerate}[label=\alph*)]
\item Jaký je sdružený výraz k $5-\sqrt2$?
\item Je pravda $(\sqrt6-\sqrt2)(\sqrt6+\sqrt2)=8$?
\item Usměrni $\dfrac{1}{\sqrt7-2}$.
\item Proč není vhodné rozšiřovat $\dfrac1{\sqrt3+\sqrt2}$ zlomkem
$\dfrac{\sqrt3+\sqrt2}{\sqrt3+\sqrt2}$?
\end{enumerate}
\tcblower
\begin{enumerate}[label=\alph*)]
\item $5+\sqrt2$;
\item Ne, výsledek je $6-2=4$;
\item $(\sqrt7+2)/(7-4)=(\sqrt7+2)/3$;
\item Ve jmenovateli vznikne druhá mocnina dvojčlenu a odmocnina nezmizí.
\end{enumerate}
\end{zadaniReseniBox}

\blok{Rozšíření: při proměnné vždy kontroluj podmínky}
\begin{prikladReseniBox}[breakable=false,title={Rozšiřující příklad 1: přípustné hodnoty proměnné}]
Usměrni $\dfrac1{\sqrt{x}-2}$ a najdi definiční obor původního výrazu.
\tcblower
Potřebujeme $x\ge0$ a $\sqrt x-2\ne0$, tedy $x\ne4$.
Pak je $\sqrt x+2>0$ a můžeme násobit:
\[
\frac1{\sqrt x-2}\cdot\frac{\sqrt x+2}{\sqrt x+2}
=\boxed{\frac{\sqrt x+2}{x-4}},
\qquad x\in\langle0;4)\cup(4;+\infty).
\]
\end{prikladReseniBox}
\begin{zadaniReseniBox}[breakable=false,title={Rozšiřující příklad 2: past s nulovým sdruženým výrazem}]
Pro $x\ge0$ se pokus usměrnit $\dfrac1{\sqrt x+1}$.
Co se stane pro $x=1$?
\tcblower
Pro $x\ne1$ lze použít
$\dfrac{\sqrt x-1}{x-1}$. V bodě $x=1$ by však
\emph{sdružený násobek} byl nula, tedy bychom násobili $0/0$.
Původní výraz je přesto v bodě $x=1$ definovaný a rovná se $1/2$.
\textbf{Ověření definičního oboru před úpravou je nezbytné.}
\end{zadaniReseniBox}

\blok{Domácí procvičování}
\begin{zadaniReseniBox}[breakable=false,title={Domácí úkol 1: sdružené výrazy}]
Usměrni a uprav:
\begin{enumerate}[label=\alph*)]
\item $\dfrac1{\sqrt5+1}$;
\item $\dfrac4{\sqrt7-\sqrt3}$;
\item $\dfrac3{4-\sqrt{13}}$;
\item $\dfrac{\sqrt3+1}{\sqrt3-1}$;
\item $\dfrac{2\sqrt2}{\sqrt6+\sqrt2}$;
\item $\dfrac{\sqrt5-\sqrt2}{\sqrt5+\sqrt2}$.
\end{enumerate}
\tcblower
\begin{enumerate}[label=\alph*)]
\item $(\sqrt5-1)/4$;
\item $\sqrt7+\sqrt3$;
\item $4+\sqrt{13}$;
\item $2+\sqrt3$;
\item $\sqrt3-1$;
\item $(7-2\sqrt{10})/3$.
\end{enumerate}
\end{zadaniReseniBox}
\poznamka{\textbf{Shrnutí:} Nejde o kouzlo: vynásobíme výrazem rovným jedné,
rozdíl čtverců odstraní odmocniny ze jmenovatele. Hodina 34 propojí
částečné odmocňování, usměrňování, mocniny a úpravy výrazů v delších úlohách.}
\end{document}
